\section{Hierarchies as Intervals}
\label{sec-labels}

An instance query $\mathrm{inst}(c)$ (Definition~\ref{def-queries}) must return
the instances of every class in $\desc(c)$.
A store can \emph{materialize} the entailed type triples, or \emph{rewrite}
the query into one scan per class in $\desc(c)$.
\sys instead numbers the classes so that $\desc(c)$ is a few integer
intervals, and keeps the numbers in the B+ tree it already has.

\stitle{Labelling}.
Collapse each cycle of $H$ into one class.
Choose a spanning forest by taking, for each class with several parents, the
parent with the longest path from a root.
Number the classes $1,\dots,|\mathcal{C}|$ in postorder of the forest, and let
$\lo(c)$ be the smallest number in the forest subtree of $c$.
In a forest, $u\in\desc(c)$ if and only if $\lo(c)\le\po(u)\le\po(c)$, a
single interval.
In a DAG each class carries a \emph{set} of intervals $I(c)$, built bottom-up
(Algorithm~\ref{alg-label}).

\begin{algorithm}[t]
\caption{Interval sets for a class DAG}\label{alg-label}
\small
\KwIn{subclass edges $E$ of $H$}
\KwOut{$\po(c)$ and an interval set $I(c)$ for each class $c$}
collapse the strongly connected components of $H$ (Tarjan), and number them parents first\;
\For{each component $c$, parents first}{
  $\mathit{depth}(c) \gets 1 + \max_{c'}\ \mathit{depth}(c')$ over its parents $c'$ \tcp*{0 for a root}
  $\mathit{tp}(c) \gets$ the parent with the largest depth\;
}
number the forest $\{(c,\mathit{tp}(c))\}$ in postorder; record $\lo(c)$\;
\For{each component $c$, children first}{
  $I(c) \gets \textsc{Merge}\big(\{[\lo(c),\po(c)]\} \cup \bigcup_{c'\text{ a child of }c} I(c')\big)$ \tcp*{sort; join overlapping or adjacent}
}
\end{algorithm}

\begin{lemma}[Correctness]\label{lem-correct}
$u\in\desc(c)$ if and only if $\po(u)$ lies in an interval of $I(c)$.
\end{lemma}
\iffull
\begin{proof}
By induction over the order in which Algorithm~\ref{alg-label} builds $I$: the
numbers covered by $I(c)$ are exactly those of the classes in $\desc(c)$.
\end{proof}
\fi

\begin{lemma}[Size]\label{lem-size}
For every class $c$, $|I(c)|\le 1+|E'(c)|$, where $E'(c)$ is the set of non-tree edges whose tail is in $\desc(c)$.
\end{lemma}
\iffull
\begin{proof}
A class $u\in\desc(c)$ is reached from $c$ by a path.
If the path uses no non-tree edge, $u$ is in the forest subtree of $c$.
Otherwise let $\epsilon$ be the last non-tree edge on it; $u$ is then in the forest
subtree of the head of $\epsilon$.
So the covered set is the forest subtree of $c$ together with the forest
subtrees of the heads of the edges of $E'(c)$, a union of $1+|E'(c)|$
intervals.
\end{proof}
\fi

The labelling takes $O(|\mathcal{C}|+|E|)$ plus the merges, and a tree needs
one interval per class.

\stitle{Index}.
The labels go into B+ trees of \sys: key $(\po(c),\mathit{id}(c),0)$ for each
class, and $(\po(c),\mathit{id}(e),0)$ for each triple $(e,\texttt{rdf:type},c)$.
The classes below $c$, or the instances of $c$ and its subclasses, are then
one range scan $[\,(a,0,0),\,(b{+}1,0,0)\,)$ per interval $[a,b]$ of $I(c)$,
and the answers of different intervals are disjoint.
Inserting an instance adds one key; deleting a class edge rebuilds the labels,
which takes tens of milliseconds for an ontology of tens of thousands of
classes.
Incremental maintenance with gaps in the numbering~\cite{agrawal1989} is not
implemented.
Dictionary ids are untouched: ids appear in three permutations, the log, and
checkpoints, and renumbering them to follow the hierarchy would rewrite all of
it.
The label tree is separate and rebuildable, which keeps ids stable.

\stitle{Membership}.
A class $u$ is below $c$ if $\po(u)$ lies in an interval of $I(c)$, a binary
search over $|I(c)|$ intervals with no scan.
This makes the class filter of Section~\ref{sec-vectors} a constant-space
predicate on a term's number.
